\documentclass[12pt]{article}

\begin{document}

An Exploration Dynamical Systems Chaos Argyris C01 Completely Revised and Enlarged

Second Edition

John Argyris, Gunter Faust

Maria Haase.Rudolf Friedrich

\hfill

\textbf{Preface to the Second Edition} 

\hfill

In 1994, when the first edition has been published in English and in German, we did
not anticipate that this introduction to the field of non-linear dynamics and chaos
would meet with such lively interest. Although attention to the subject has lessened
in the media and in popular scientific publications over the last 20 years, the
fundamental ideas, the theoretical insights and the tools developed for the analysis
of non-linear dynamical systems have continued to spread into the most diverse
areas of science and technology as well as into the respective literature and are
nowadays part of the classical curriculum in many study programmes. Since the
English edition sold out many years ago, we decided to follow up the expanded and
revised second German edition with another extended edition in English.

\hfill

After John Argyris passed away in 2004 at the advanced age of 91, we considered it
our obligation to continue work on this book in the original sense – first in a second
German edition – attempting to explain the complex topic of non-linear dynamics
to a wide audience as descriptively and vividly as possible.

\hfill

We took the opportunity to emphasise and further clarify the differences between
purely temporal and spatio-temporal dynamics, in particular between chaos and
turbulence. The initial hope – occasionally accompanied by euphoria – that insight
into the special features of chaotic systems meant that scientists were about to
solve the centuries-old problem of turbulence was not fulfilled. Without doubt, the
elucidation of characteristics such as unpredictability and mixing in chaotic systems
made a substantial contribution to understanding turbulent flows, yet the fundamental
questions concerning fully developed turbulence are unsolved to this day.

\hfill

By enlarging the circle of authors to include Rudolf Friedrich, we were able to add
with Chapter 9 a completely new and more ambitious chapter on turbulence, which
comprises and critically scrutinises the fundamentals of turbulent flows and a series
of classic turbulence models. While there is a multitude of excellent monographs on
turbulence from the viewpoint of the engineering sciences, the new chapter mainly
discusses the fundamental questions from the physical point of view, highlighting
the common ground with concepts existing in chaos theory. To describe and investigate
fully developed turbulence, it is indispensable to include probabilistic and
stochastic methods. So that the reader can understand the chapter on turbulence
better, we thought it appropriate to include two new sections in the mathematical
introduction in Chapter 3: basic concepts of probability theory and invariant measure
and ergodic orbits. Our aim in these sections was to present the basic ideas and
concepts; we did not, however, strive for a complete description, rather referring
the reader to the excellent literature on probability and stochastic processes.

\hfill

In section 3.8.5, we added a brief introduction to wavelet transformation and outlined
its relevance for time-frequency analyses of time series. The use of wavelets
is especially favourable for analysing multifractal structures as described in section
8.5.2. In section 3.9.12, we also included a brief description of Markov analysis,
a recently developed and far-reaching method which allows us to separate the deterministic
part of the dynamics from the dynamical and even measurement noise. In
section 9.6.8, this method is applied to velocity increments measured in turbulent
flows. Section 8.7 contains an updated and extended section on routes out of chaos.
Shilnikov bifurcations and spiral chaos received much attention in various applications,
from chemical reactions to the propagation of nerve impulses and epileptic
seizures. To understand such systems better, we also included a short overview in
section 10.6. In addition, the paragraph on the kinetics of chemical reactions on
surfaces has been largely revised in the hope of rendering this widely used field of
application more comprehensible to a larger audience of students, see section 10.9.
On the other hand, we dispensed with the section on celestial mechanics; this would
have gone beyond the scope of this monograph.

\hfill

Suddenly and unexpectedly, on 16th August 2012, Rudolf Friedrich passed away,
taken from the midst of his life and work. His broad knowledge and interdisciplinary
expertise made him extraordinarily inventive, not only in theoretical physics, his
specialist field, but also in many other areas of science, including experimental and
technical applications. We hope we have been able to complete the work on this
book as he would have wished.

\hfill

Without the vigorous support of numerous friends and colleagues from many scientific
disciplines, we would not have been able to finish this extended and updated
second edition. To implement the original KeTEXversion in LaTeX and to generate
and reconstruct the figures with the programme system AnT 4.669, we received professional
help. We acknowledge the valuable hints for additions and improvements
which we obtained in the course of many discussions and during proof-reading.
We wish to thank all those who constantly and sedulously helped us and are especially
indebted to Viktor Avrutin, Inna Avrutina, Rolf Bader, Anton Daitche,
Martin Dziobek, Markus Eiswirth, Jan Friedrich, Jason Gallas, Svetlana Gurevich,
Andreas Haase, Marion Hackenberg, Marcus Hauser, Oliver Kamps, David
Kleinhans, Yuri A. Kuznetsov, Bernd Lehle, Pedro Lind, Johannes L¨ulff, Joachim
Peinke, Peter Plath, G¨unter Radons, Michael Schanz, Susanne Schmidt, Daniel
Stellbrink, Robert Stresing, Daniel H. Sugondo, Christian Uhl, Hans van den Berg,
Judith Vogelsang, Georg Wackenhut and Michael Wilczek. Our special thanks go
to Prudence Lawday, the translator of the first English edition, for her excellent
help with this extended manuscript. We are also indebted to Michael Resch, the
director of the High Performance Computing Center Stuttgart (HLRS) of the University
of Stuttgart, for his generous assistance and to the Argyris Foundation for its
financial support.

\hfill

Last but not least, we wish to thank the members of staff at Springer-Verlag for
their professional co-operation, and in particular Thomas Ditzinger for his patience
and continuous encouragement and support. Our sincere thanks go to everyone who
supported us.

\hfill

Stuttgart, October 2014 M. Haase, G. Faust

\hfill

\textbf{1 Descriptive Synopsis of the Text}

But thought’s the slave of life, and life time’s fool;

And time, that takes survey of all the world,

Must have a stop.

William Shakespeare, King Henry IV, Part I, IV, 2, 81

\hfill

This book is conceived as an elementary introduction to the modern theory of nonlinear
dynamical systems with particular emphasis on the exploration of chaotic
phenomena. One might ask why yet another book should be published when the
literature on chaos and non-linear oscillations already fills shelf after shelf following
the stormy developments in this branch of science since the 1970s. The reasons
which prompted us have been detailed in the preface.

\hfill

Our aim was to write an elementary, yet detailed text-book, without assuming an
all-too profound mathematical knowledge on the part of the reader. It is aimed
not at mathematicians and theoretical physicists, but rather at budding scientists
such as chemists, biologists and engineers in all fields, particularly in the aeronautical
and electrical engineering, but also in musical acoustics, as well as economic
scientists and interested graduate students. Should they be required for a comprehension
of the subject, mathematical procedures and tools are introduced in small
sub-sections; in all cases, we have attempted to support the text illustratively by
graphics, colour plates and additional computer demonstrations. A fairly extensive
list of references is intended to help the reader in his quest for further reading.

\hfill

Chapters 2 and 3 are an elementary introduction to the intellectual and conceptual
edifice of non-linear dynamical systems. In Chapter 2, we present an outline
of the complex of problems investigated in this book. We consider non-linear dynamical
processes, restricting ourselves to deterministic and, with few exceptions,
finite-dimensional systems which can be described by systems of ordinary differential
equations. In contrast to classic treatises on the theory of oscillations which generally
devote a great deal of space to the solution of linear problems and mainly consider
special solutions of selected problems, the theory of dynamical systems founded by
Henri Poincar´e allows qualitative statements on the overall set of the solutions of
specific as well as broad classes of differential equations by a combination of analytical
and geometrical, differential topological methods. Central questions concern
the structure of the phase space, the stability behaviour of the solutions, their longterm
response and the effects of bifurcations, i.e. the qualitative changes in the flow
as a result of variations in the parameters. The fact that even the simplest nonlinear
systems can generate chaotic behaviour leads to the discussion in section 2.1 of
the concept of causality and the connection between determinism and predictability.
Subsequently, using the example of single-degree-of- freedom oscillators with and
without friction, with and without excitation force, we describe possible states of
motion of dynamical systems – such as stationary, periodic and chaotic behaviour
– and introduce some elementary basic concepts such as phase space, trajectories,
manifold, attractors etc.

\hfill

A central theme in the analysis of non-linear systems is the discussion of the stability
characteristics of states of equilibrium, periodic motions etc. Chapter 3 thus
begins with a discussion of the singularities of linear differential equations and a
classification of the fixed points on the basis of eigenvalues and eigenvectors for the
two-dimensional case as well as the possibilities of an extension of the results to
non-linear systems. The introduction of a series of stability concepts is followed by
the definition of the invariant manifolds which enable a structuring of the phase
space. We then discuss Poincar´e sections and discrete maps, followed by a first
elementary meeting with the logistic map where we can observe a transition from
stationary to periodic, multiple periodic behaviour to chaos through variations of
a control parameter. General periodic motions can be easily described in Fourier
space. Therefore, we continue with a brief introduction to harmonic analysis, including
a short survey of the most important properties of the Fourier transform
and a reference to wavelet analysis, which allows us to perform time-frequency analyses
of non-stationary processes.

\hfill

The subsequent section 3.9 presents a summary of the basic principles of the theory
of probability. Although this book is devoted to the analysis of deterministic nonlinear
systems, one often needs a set of fundamental terms and definitions stemming
from probability theory to characterise the dynamics of chaotic systems in
phase space and to describe turbulence. Among these are concepts such as random
variables, probability density, conditional and joint probability and moments of distribution
functions. We briefly outline the importance of the central limit theorem
and present a short description of the most commonly used density distributions.
Finally, we give a short overview of the recently developed Markov analysis, which
allows for sequentially measured experimental data to disentangle the underlying
deterministic dynamics from the often unavoidable superimposed fluctuating forces
and to extract the basic stochastic equations. On the basis of probability theory, we
introduce in the final section 3.10 the concepts of invariant measure and ergodicity.

\hfill

Dynamical systems can be assigned to two basic families, namely conservative systems
not exposed to loss of energy and dissipative systems subject to loss of energy.
In Chapter 4, we consider the conservative systems and attempt to span the gap
between integrability and chaos. Within the concept of classical mechanics, the
Hamilton equations are derived, taking d’Alembert’s principle as a starting point.
The theorem of Liouville allows us to investigate the volume conservation of the
phase flow of conservative systems and thus to distinguish between conservative
and dissipative systems. An analytical solution of non-linear Hamilton differential
equations is only possible in exceptional cases, namely if the system can be transformed
to action and angle variables and thus described by quasi-periodic motions
on a torus in phase space. Such an integration is demonstrated on the example of
Kepler’s motion of a planet around the sun. In contrast, an analytical solution of
the three-body problem is no longer possible. Towards the end of the 19th century,
Poincar´e succeeded in proving that small perturbations of the two-body problem
throw doubt on the stability of the solutions. He showed that perturbation methods
do not necessarily lead to correct results since the series used in the course
of the calculation might diverge, thus possibly leading to false conclusions. It was
not until 60 years later that, within the framework of the KAM theory, Andrei
N. Kolmogorov, Vladimir I. Arnold and J¨urgen Moser were able – under certain
conditions – to answer the question of the stability of perturbed Hamilton systems.
In section 4.4, the KAM theory is outlined and the disintegration of the tori in phase
space demonstrated on Poincar´e sections of two-dimensional perturbed Hamilton
systems. It turns out that the emergence of heteroclinic points generates the irregular,
chaotic behaviour in the neighbourhood of hyperbolic points. A detailed
discussion of the conservative H´enon map closes this chapter.

\hfill

Chapter 5 deals with dissipative systems. In the classical mechanics of Galilei and
Newton, loss of energy due to friction was considered a disturbing factor and was
thus neglected. However, this idealisation can be applied in a good approximation
to the orbits of the planets, but only rarely to motions on earth. The theory of
dynamical systems shows that the introduction of dissipative terms does not make
the problem more complex, but rather simplifies the long-range solution in many
cases. The long-term dynamics of high-dimensional dissipative systems is often conditioned
by some few fundamental modes; a contraction of the phase space volume
takes place and attractors emerge which cannot be generated in conservative systems.
Dissipative systems not subject to energy supply come to rest, they approach
a state of equilibrium or point attractor. If the system is fed with energy, then, in
the simplest case, periodic or quasi-periodic behaviour can occur or, on the other
hand, completely irregular behaviour which is described in phase space by a strange
attractor, an extremely complex structure with zero volume. The typical characteristics
of such strange attractors are summarised using the example of the Lorenz
attractor, the structure of which is characterised fundamentally by the sensitive
dependence of the trajectories on the smallest perturbations of the initial conditions.
In the subsequent sections of Chapter 5, we review a series of mathematical
tools – namely the power spectrum, autocorrelation, the Lyapunov exponents and
the dimensions in phase space – which allow a quantitative characterisation of and
differentiation between the individual attractor types.

\hfill

We begin with two classic criteria, the power spectrum, which supplies information
on which basic frequencies dominate, and autocorrelation, which reflects the memory,
i.e. the temporal interrelations of a signal. Both criteria are based on Fourier
transformation, the most important characteristics of which are summarised previously
in section 3.8. While demonstrating a differentiation between regular and
irregular behaviour, they are not sufficient to characterise chaotic motion. Further
mathematical tools are necessary, such as the Lyapunov exponents, for example,
which provide a criterion for the stability behaviour of neighbouring trajectories.
The classic Floquet theory, which allows a linear stability analysis of periodic motions,
is extended to irregular courses of motion and leads to the spectrum of the
Lyapunov exponents. Section 5.4 closes with the description of algorithms which
allow a numerical calculation of all Lyapunov exponents.

\hfill


The concept of dimension represents a further important instrument for the classification
of the various attractor types. Damping generally ensures that the number
of independent variables determining the dimension of the phase space are reduced
considerably during the transient phase. The dimension of an attractor is thus a
suitable source of information on the number of fundamental modes participating
in the long-term response: a point attractor of dimension 0 corresponds to a state
of equilibrium, a limit cycle of dimension 1 to a periodic motion etc. A strange
attractor, which, for example, is a sort of hybrid between a surface-like and a spatial
structure in a three-dimensional phase space, corresponds to a chaotic motion.
Classic dimension concepts are not adequate to describe its intricate, mille-feuillelike
structure; it has to be ascribed a non-integral dimension between 2 and 3.

The French mathematician Benoˆıt Mandelbrot coined the word “fractals” for such
sets. What is responsible for the complexity of fractal sets is their self-similarity,
i.e. their invariance with respect to alterations in scale. There is a whole family of
suitable dimension concepts – such as the capacity dimension, the information and
correlation dimension etc. – which we shall discuss in detail in section 5.5.

\hfill

Dimensional concepts are often used to characterise experimental time series. If a
system is based on deterministic laws, the temporal evolution of a single variable
contains all the information on the total dynamics of the system. We discuss a
reconstruction technique which offers the opportunity of reconstructing attractors
from one-dimensional time series and also applying in this context concepts such
as dimension (section 5.5.4) and the Lyapunov exponents (see also section 10.9) to
characterise the associated orbits.We conclude Chapter 5 with the discussion of the
Kolmogorov-Sinai entropy, based on Shannon’s information theory and quantifying
the degree of disorder of a system. With the example of a simple map, the Bernoulli
shift, it is demonstrated that the KS entropy measures the information transport
from a microscopic to a macroscopic scale. Moreover, it allows for the specification
of a bound for the time-span within which reliable prognoses on the course of the
trajectories are still possible.

\hfill

Typically, one or more control parameters appear in the model equations describing
physical systems. If these differential equations have non-linear character, a variation
of the parameters can result in a qualitative change of the topological structure
of the solutions. Such changes are referred to as bifurcations. Chapter 6 contains a
detailed discussion of the methods of the local bifurcation theory; we concentrate
on local bifurcations of states of equilibrium and periodic motions arising from the
variation of one control parameter. As an introduction, we first present some typical
bifurcations of fixed points and elucidate the qualitative change of the solutions
with corresponding phase portraits.

\hfill

In the neighbourhood of a hyperbolic fixed point, it suffices to investigate the linearised
system in order to establish the stability and solution behaviour. At bifurcation
points, however, a fixed point loses its hyperbolic character. In this case,
non-linear terms must be included in our considerations in order to clarify the stability
and the solution behaviour. To this end, the method of the centre manifold is
applied. This approximation technique, presented in section 6.2, leads to a drastic
reduction of the number of equations necessary to describe the long-term behaviour
in the neighbourhood of non-hyperbolic fixed points and hence the clarification of
the stability.

\hfill

It is the purpose of the local bifurcation theory to classify the local bifurcations of
fixed points, limit cycles etc. and to set up their normal forms. In section 6.3, we
describe the method of normal forms, a technique with which non-linear differential
equations in the neighbourhood of fixed points and limit cycles can be reduced to
the simplest form by a series of non-linear transformations; the non-linear terms are
eliminated successively whenever possible. A complete linearisation is only possible
if no resonances occur. Normal forms of bifurcations, however, necessarily contain
resonances.

\hfill

In section 6.4, we discuss the bifurcations of one-parametric flows, deduce the corresponding
normal forms and list the conditions which allow a classification of the
individual bifurcation patterns. We then consider the question of the robustness
of the individual bifurcations, i.e. we investigate whether small alterations in the
model equations generate qualitative changes of the bifurcation patterns. Finally,
we discuss the normal forms of bifurcations of fixed points of one-parametric maps.
Since periodic orbits appear as fixed points in the Poincar´e map, the discussion also
includes the bifurcations of limit cycles.

\hfill

In section 6.7, we demonstrate with the example of the logistic map that a continuous
increase of the control parameter can generate a series of local bifurcations,
in this case a cascade of period doublings, ultimately leading to chaotic behaviour;
this latter is interrupted repeatedly by windows of regular behaviour, however.
Mitchell Feigenbaum was the first to prove that this route to chaos is not restricted
to the special form of the logistic map, but rather is a universal scenario which
can be observed in many systems. In this section, we explore the background to
the mechanism of continued period doublings and explain how the self-similarity
characteristics of the bifurcation diagram can be decoded with the aid of renormalisation
techniques and how the Feigenbaum constants can be determined. We then
turn to the link with second-order phase transitions, known to us, for example, from
the transition from a liquid to a gaseous state or from the ferromagnetic transition.

\hfill

We conclude Chapter 6 with an introductory outline of synergetics, a branch of
science founded at the end of the 1960s by the physicist Hermann Haken which
investigates the dynamics of pattern formation processes. The central question of
synergetics is whether and to what extent the formation of structures in the most
varied realms of physics, biology, sociology, medicine etc. in animate and inanimate
nature is based on generally valid principles and how this self-organisation can be
reproduced in a mathematical formalism. On the one hand, we refer to the features
in common with the centre manifold theory; on the other, we draw attention to
the more extensive framework of synergetics, a result of the inclusion of spatial
inhomogeneities and fluctuating forces.

\hfill

Chapter 7 is devoted to the structure formation in convection flows and the deduction
and discussion of the Lorenz system, which stimulated the development of
the chaos theory fundamentally. The chapter begins with a qualitative description
of the physical laws leading to the formation of the characteristic convection rolls
in the Rayleigh-B´enard experiment. There follows a detailed deduction of the basic
hydrodynamic equations for convection problems. The continuity equations, the
Navier-Stokes equations and the heat equation are deduced from the conservation
theorems of mass, impulse and energy. In many cases, these basic equations can
be simplified by means of the Boussinesq-Oberbeck approximation, in which the
variation of the density is only accounted for in the buoyancy term.

\hfill

In the early 1960s, the American meteorologist Edward N. Lorenz applied this system
of basic equations to study the accuracy of weather forecasts. While his highly
simplified model for the B´enard problem cannot stand up to an experimental verification
describing the onset of turbulence, it nevertheless led to a new paradigm in
the natural sciences: On the one hand, there are basic limits to the predictability of
a dynamical system in the latent presence of chaotic behaviour; on the other hand,
however, time series which have been classified as purely stochastic in the past have
turned out to describe in certain cases chaotic motions, the essential definition of
which can often be construed with a restricted number of non-linear equations.

\hfill

In section 7.3, we follow the original publications of Saltzman and Lorenz and deduce
the three ordinary differential equations of the Lorenz system. Subsequently,
the evolution of the solutions of this system is described for variations of the free
parameter which controls the temperature difference between the upper and the
lower plate in the Rayleigh-B´enard problem; the emergence of local and global bifurcations
is observed. A series of colour plates demonstrates on the one hand the
spatial and temporal evolution of the convection rolls in the Lorenz model and
on the other hand the significant changes of the phase space structure due to an
increase of the control parameter; these are illustrated by consideration of the evolution
of the stable manifolds.

\hfill

In Chapter 8, we present a series of mathematical models describing the transition
from regular to chaotic motions. In the early years of chaos research, the hope was
that these scenarios could describe the transition to turbulent behaviour, thus making
a substantial contribution to the solution of the turbulence problem.

\hfill

The first model describing the onset of turbulence goes back to Lev D. Landau
(1944), who assumed that the transition from laminar to turbulent behaviour is
generated by an endless series of (Hopf) bifurcations. It was not until a quarter of
a century later that the deficiencies of this model were recognised: Landau’s model
cannot explain the sensitive dependence of the flow on small perturbations nor generate
a mixing of the trajectories in the phase space. In 1971, David Ruelle and
Floris Takens proved that generally quasi-periodic motions on higher-dimensional
tori are not stable, as Landau had presumed, but rather that the quasi-periodic
motion becomes chaotic abruptly after the emergence of the third incommensurable
frequency, i.e. the motion on a T 3-torus collapses into a strange attractor,
a designation coined by Ruelle and Takens. A great step forward had been taken
and a window opened towards the description of the characteristics of turbulence
– such as unpredictability of the motion and mixing – on the basis of deterministic
equations. With the aid of refined measuring techniques, Gollub and Swinney
succeeded in achieving an experimental verification for the B´enard problem based
on the analysis of the power spectra and supported by the determination of the
dimension of reconstructed attractors.

\hfill

Sections 8.3 to 8.5 contain a discussion of the transition from quasi-periodic to
chaotic motions using the circle map as a model. Two different scenarios are conceivable.
Either chaos occurs following a synchronisation of the frequencies and a
subsequent series of period doublings – this is the Feigenbaum route already discussed
in section 6.7 for the logistic map – or there is a direct transition from
quasi-periodicity to chaos. The second route requires the tuning of two control parameters
and can also be interpreted by means of the circle map. Both transitions
to chaos are of universal character and the scaling characteristics can be decoded
in both cases with the aid of the renormalisation theory. A selection of figures aims
at facilitating access to this complex material. There follows a description of the
work of Albert Libchaber who succeeded in confirming the theoretical prognoses
quantitatively by experiments which were highly precise at that time. Section 8.6
contains a discussion of intermittent transitions which, following Yves Pomeau and
Paul Manneville (1980), can be classified in three groups.

\hfill

Having presented a series of scenarios describing transitions from regular to chaotic
behaviour, we present in section 8.7 an overview of various strategies allowing the
reverse, a conversion of chaotic motions to regular regions. Following the seminal
paper of Edward Ott, Celso Grebogi and James A. Yorke from 1990, this issue became
a field of non-linear dynamics which experienced a rapid development. This is
because it was recognised in many areas like medicine, chemistry, laser technology
and telecommunication, as well as in other fields of engineering science, that chaotic
behaviour contains a high potential with regard to a possible control and regulation
of chaos. To this end, one takes advantage of a characteristic feature of chaos, the
sensitive dependence of the dynamics on small perturbations, by imposing minimal
control impulses which stabilise the unstable orbits inherently contained in strange
attractors. In section 8.7 we give a brief survey of the most effective strategies and
provide information on a multitude of applications with the potential to lead to
innovative developments.

\hfill

In the beginning, it was hoped that chaos research would help uncover the basic
mechanisms leading to fully developed turbulence. These high expectations have
not yet been fulfilled. The reason is that chaos theory researchers initially confined
themselves to a purely temporal description of the dynamics, which is not sufficient
for an understanding of turbulent flow. It is rather the additional spatial dependence
of all quantities which leads to complex interrelations and to an extremely
high dimensionality of the corresponding phase space.

\hfill

Chapter 9 presents an overview of the state of turbulence research from the present
point of view, almost 40 years after the paradigm shift in the understanding of
non-linear processes. On the one hand, we identify the issues in which chaos theory
has made substantial contributions to an understanding of turbulent flow. On
the other hand, it becomes evident that the spatio-temporal complexity and the
extraordinarily high number of excited modes in fully developed turbulence makes
the use of statistical methods unavoidable.

\hfill

In the chapter on turbulence, the main focus is on various vortex models and the
discussion of their influence on the dynamics of fluid flow. A detailed description of
the characteristic scales, of the energy cascade and of Heisenberg’s turbulence model
follows. Particular emphasis is put on the derivation of the K´arm´an-Howarth equation
and an exposition of the underlying assumptions. Subsequently, we describe
the classical turbulence models of Kolmogorov from 1941 and 1962 as well as more
recent multifractal models presented by Frisch and She-Leveque and discuss the
strengths and weaknesses of the respective models. We also report on an application
of the Markov analysis presented in section 3.9 to the velocity increments of
a turbulent jet. The fact that the turbulent data can be described by a Markovian
process in scale allows us to express the full n-point statistics by the 2-point transition
probability densities and to reproduce the intermittent behaviour accurately.
Finally, we add an overview of some open questions and possible issues relevant for
future research.

\hfill

In the final Chapter 10, we present some fields in which the application of the theory
of dynamical systems is very promising. With the aid of computer experiments
we discuss a variety of classical models displaying features which are characteristic
of a multitude of dynamical systems.

\hfill

For non-linear systems, analytical solutions can only be given in exceptional cases;
in the case of chaotic behaviour, we generally have to rely completely on numerical
investigations. This is the explanation for the fact that, as a result of efficient computers,
the theory of dynamical systems – once the domain of the mathematical and
physical faculties – is now receiving increasing notice in the most varied disciplines
such as biology, chemistry, meteorology, electrical engineering and hydrodynamics,
as well as in economic sciences, as demonstrated by the flood of publications from
these fields.

\hfill

The discovery of universal characteristics, i.e. of fundamental principles common
to all non-linear systems of a whole class of problems, was without doubt a milestone
in the more recent history of the theory of dynamical systems and justified
the intensive study of seemingly over-simplified systems. The investigations were
accelerated greatly by the swift development of efficient computers – indeed, they
were only made possible through this explosive growth of computer power.

\hfill

The first section of Chapter 10 offers an insight into bone remodelling processes. In
order to achieve a durable anchoring of endoprostheses, models are being developed
with the aim of describing bone retrogression or bone formation on the basis of dynamical
systems. Should this succeed, computer simulations could speed up the
development of more reliable types of endoprosthesis. In a first study, two very simple
dynamical systems for the modelling of bone adaptation processes are discussed.
Specially constructed recursion rules are particularly suitable for elucidating general
characteristics of dynamical systems due to the low computational effort. In
particular, the dissipative H´enon map was conceived as a model for a Poincar´e
map of a three-dimensional system. With the aid of this map, we demonstrate in
section 10.2 that a mixing of the trajectories in the phase space results from a continuous
repetition of stretching and folding processes. This mechanism is the key
to the understanding of the exponential drifting apart of neighbouring trajectories
in the phase space; it is illustrated with the example of the Lorenz attractor in a
series of colour plates.

\hfill

In section 10.3, we return to the Lorenz equations and discuss some further characteristics
of the system. A number of figures demonstrate the emergence of cascades
of period doublings within regular regions and the coexistence of attractors. The
bifurcation diagram of the Poincar´e sections is highly reminiscent of that of the logistic
map, thus offering evidence of universal laws. Conclusions about the system
behaviour can be drawn by calculating the Lyapunov exponents. The section ends
with the determination of the capacity dimension of the Lorenz attractor according
to Hunt and Sullivan’s procedure.

\hfill

In section 10.4, we turn to the van der Pol equation, originally set up to describe
electrical circuits. The system possesses a non-linear damping term with the
characteristic that the system is supplied with energy for small amplitudes, which
manifests itself in a negative damping term, whereas energy is dissipated for large
amplitudes. Typically, self-excited oscillations occur in such systems, i.e. although
no external periodic forces are applied, the system tends towards a limit cycle on
which energy supply and dissipation are in balance in the temporal mean. We first
consider the van der Pol equation without external excitation. In the case of a
small damping parameter, i.e. weak non-linearity, we observe, as expected, almost
sinusoidal oscillations; in the case of large damping values, on the other hand, we
observe self-excited relaxation oscillations. Computer investigations suggest that
exactly one limit cycle exists in the phase space; the strict mathematical proof of
this observation with the aid of the Poincar´e-Bendixson theorem is laborious. If
small periodic excitation forces are applied in the case of low damping, the averaging
method can be used; here, by taking a time average, the initial system is
replaced by an autonomous system involving the neglect of higher harmonics. Large
damping parameters respectively strong non-linearities together with periodic external
excitation are discussed for a modification of the van der Pol equation. On
the basis of Poincar´e sections and phase portraits, the character of the chaotic solutions
on the Birkhoff-Shaw attractor is illustrated.

\hfill

The Duffing equation considered in the subsequent section represents one of the
simplest dynamical systems in which chaotic behaviour can be observed. Differential
equations of its type are obtained in the formulation of non-linear pendulum
oscillations or for electrical circuits and other problems, for example. It approximates
the behaviour of an externally excited non-linear oscillator with one degree of
freedom. We deduce the Duffing equation for the example of an externally excited
buckled beam. In this case, the equation of motion is a non-linear partial differential
equation, the solution of which is approximated with the aid of the Galerkin procedure
taking the first fundamental mode into account. As a result, we obtain the
Duffing equation with negative stiffness. According to the selection of the control
parameters in the Duffing equation, we observe a great variety of system response.
The Duffing equation is integrated numerically for a series of parameter values and
the results presented in the form of phase portraits and Poincar´e sections.

\hfill


A surprising phenomenon can emerge in dissipative systems when several attractors
coexist in the phase space. Even for regular behaviour, the boundaries between the
basins of attraction can assume fractal character so that predictions on the longterm
behaviour are no longer possible in this domain. With the aid of the Melnikov
method, bounds for the control parameters for which fractal boundaries emerge can
be calculated explicitly. For the example of the Duffing equation, the calculation of
the so-called Holmes-Melnikov boundary is carried out explicitly.

\hfill

In section 10.6, we briefly discuss some bifurcations of flow involving homoclinic
orbits in three-dimensional vector fields. The most complex case, the so-called
Shilnikov bifurcation, leads to a highly complex topological structure of the flow
and can be observed, for example, in the propagation of nerve impulses along axons
or during epileptical seizures.

\hfill

By now, computer experiments play an important role in non-linear dynamics. In
general, theorems cannot be proven with the aid of computers; nevertheless, numerical
investigations in association with modern computer graphics reveal unexpected
characteristics and links in non-linear systems which had previously not been taken
into account. Seen in this way, numerical experiments are valuable aids nowadays.
The two sections 10.7 and 10.8 illustrate the fact that graphic representations can
be extraordinarily helpful for the comprehension of some phenomena occurring in
non-linear recursion rules.

\hfill

Shortly after the FirstWorldWar, the two French mathematicians Gaston Julia and
Pierre Fatou studied the iteration of rational functions. They discovered that, depending
on the initial value, the iteration sequences converged to different solutions;
the boundaries between the basins of attraction of these attractors exhibit highly
peculiar characteristics and are of unsuspected multifariousness. Due to the unavailability
of computers, what was disclosed to the inner eye of Julia or Fatou remained
hidden to the average scientist until, around 1980, Mandelbrot again turned to one
of the simplest non-linear iteration rules in the complex plane, exploiting the opportunities
offered by high-performance computers. Section 10.7 offers an insight into
the manifold structures of Julia sets and the fractal boundaries of the Mandelbrot
set. Calculation of the basin boundary of the point attractor $z= \infty$ for constant values
of the control parameter of the mapping rule yield, depending on the parameter
value, completely differing limit sets, the so-called Julia sets.What is fascinating
about Mandelbrot’s recursion formula is that, in spite of its simplicity, such structural
diversity emerges. A series of computer graphics shows connected fractal sets,
dendritic forms as well as burst sets. The Mandelbrot set, on the other hand, encompasses
all values of the complex control parameter for which the Julia sets are
connected. A series of colour plates presenting various selective enlargements depicts
the richness of structures in the fractal boundaries of the Mandelbrot set.
In section 10.8, we again turn to the circle map of Chapter 8 and, by means of the
Lyapunov exponents, gain an overview of the inner structure of the Arnold tongues.
The results agree very well with the theoretical results. On the basis of the Lyapunov
exponents, the different transitions to chaos can be discussed very simply.

\hfill

Their self-similar structure is reflected in the corresponding bifurcation diagrams.
In the last section of Chapter 10, section 10.9, we decided to present an overview of
some interesting dynamic manifestations from the field of physical chemistry which
display chaotic outbursts.We concentrate on the dynamic phenomena arising in the
oxidation of CO on the surface of platinum crystals. Significantly enough, detailed
experiments have been performed under isothermal low-pressure conditions. Extensive
theoretical and experimental research work of high originality and distinction
has been performed on this subject at the Fritz Haber Institute of the Max-Planck-
Gesellschaft in Berlin under the direction of Gerhard Ertl.

\hfill

In these extensive and varied numerical investigations and parallel experiments, a
broad spectrum of complex temporal behaviour was observed, such as mixed-mode
oscillations, a transition via period doublings to chaos and even hyperchaos, characterised
by at least two positive Lyapunov exponents.

\hfill

All these phenomena represent purely temporal evolutions as they would arise, for
example, in a well-stirred medium; note, however, that in a heterogeneous reaction
on a surface, stirring is not feasible. As a result, in heterogeneous reactions, more
complex phenomena involving spatial dependencies will arise and complicate the
model and the numerical analysis of such reactions considerably.

\hfill

In a survey on CO oxidation on platinum crystals, Georg Ertl also discussed extensive
experimental studies on spatio-temporal phenomena observed by him and his
collaborators at the aforementioned Fritz Haber Institute. The formation of stationary,
propagating as well as soliton waves and even spiral chemical waves was
recorded. Irregular behaviour, such as that of chemical turbulence, also emerged.

\hfill

All the numerical investigations of the Fritz Haber Institute are based on kinetic
models evolved for the description of the temporal dynamics on a platinum (110)
single crystal surface. The rate constants and parameters arising in the equations
have been studied extensively using a variety of surface science techniques in conjunction
with complex detection and recording devices. We mention only parenthetically
the application of low-energy electronic diffraction (LEED) studies of
oscillations on platinum crystals, based on the motion of slow electrons; here, the
researchers also observed periodic changes on the surface structure which oscillated
between the (1Å~1) and the respective reconstructed (1Å~2) phase. Such structural
surface changes in turn influence the evolution of other elementary processes and
demand additional refinements of the model.

\hfill

In view of the increasing complexity, the researchers at the Fritz Haber Institute
refined the analytical model of oxidation stepwise, starting with the simplest simulation
and then refining it step by step to account for additional effects and variables
in order to describe the dynamical behaviour of increasing complexity of the CO
oxidation on platinum (110) crystals. We elaborate on the importance of the individual
steps since they clarify in an examplary manner the stepwise development of
models by illustrating the extent of validity and the limits of the respective models.
We close this section with a short look at recent applications in biochemistry. Biochemical
reactions involving enzymes which play the role of catalysts are, of course,
much more multifarious and complex than the previously mentioned catalytical
procedures. However, since biological systems are open systems exchanging energy
and mass with their environment and since the occurring processes have a nonlinear
character, it is once again indispensable to apply methods from non-linear
dynamics and synergetics in order to take at least a small step towards understanding
the complex interrelations of the system.


\end{document}


